\documentclass{article}
\usepackage[margin=1in]{geometry}
\usepackage{graphicx}
\usepackage{amsmath}
\usepackage{booktabs}
\usepackage{placeins}
\usepackage[hypertexnames=false]{hyperref}
\renewcommand{\figurename}{Fig.}

\graphicspath{{plots/}}

\title{Cluster profile emulation:\\
CGD, CGED, CEP, CEEP, CMP, CPP, CTP, CYP}
\author{}
\date{}

\begin{document}
\maketitle

\section{Profile set and conventions}

We emulate eight radial cluster profiles as functions of the scaled
radius $r/R_{500c}$, on a common $80$-bin grid restricted to
$r/R_{500c}\in[0.015,2.75]$:

\begin{center}
\begin{tabular}{lll}
\toprule
Tag & Quantity & Label \\
\midrule
CGD  & Cluster gas density & $\rho_{\rm gas}/\rho_{\rm crit}$ \\
CGED & Cluster gas electron density & $n_e$ [cm$^{-3}$] \\
CEP  & Cluster gas entropy & $K$ [keV cm$^2$] \\
CEEP & Cluster electron entropy & $K_e$ [keV cm$^2$] \\
CMP  & Cluster gas metallicity & $Z/Z_\odot$ \\
CPP  & Cluster gas pressure & $P$ [keV cm$^{-3}$] \\
CTP  & Cluster gas temperature & $T$ [keV] \\
CYP  & Cluster Compton-$y$ (tSZ) & $y$ \\
\bottomrule
\end{tabular}
\end{center}

Multi-snapshot emulators are trained per profile over the five
snapshots with clean profile data ($z\in[0,0.5]$, snapshot index 6--10),
using SEPIA's PCA basis with \texttt{exp\_variance}=0.999. The five
held-out test sims are $\{3,11,19,27,35\}$.

\section{Profile ensemble across redshift}

\begin{figure}[htbp]
    \centering
    \includegraphics[width=0.9\linewidth]{profiles_all_multiz.png}
    \caption{Training ensemble of all eight cluster profiles at three
    representative redshifts (columns) within $z\in[0,0.5]$. Each black
    line is one simulation. Profiles rotate from gas density at the top
    to Compton-$y$ at the bottom; the radial dependence and amplitude
    spread is visible in every observable. Inner-cluster behaviour
    ($r<0.1\,R_{500c}$) shows the largest sim-to-sim variation,
    consistent with feedback being strongest in the core.}
    \label{fig:profiles_all_multiz}
\end{figure}

\FloatBarrier

\section{Hold-out validation at $z=0$}

\begin{figure}[htbp]
    \centering
    \includegraphics[width=1\linewidth]{profiles_all_validation_z0.png}
    \caption{Hold-out validation at $z=0$ for all eight profiles.
    Each panel overlays three test simulations (solid) and the emulator
    prediction (dashed) with a $90\%$ Gaussian band. The emulator
    tracks the simulation shape across the full radius range; residual
    misfits at the outermost bin are dominated by low-cluster-count
    statistics in the simulation.}
    \label{fig:profiles_valid_z0}
\end{figure}

\begin{figure}[htbp]
    \centering
    \includegraphics[width=1\linewidth]{profiles_multiz_validation.png}
    \caption{Multi-$z$ validation grid across all eight profile
    emulators, evaluated on three test sims at every trained snapshot.
    Black: simulation. Red: emulator prediction with $90\%$ band.}
    \label{fig:profiles_multiz_valid}
\end{figure}

\subsection{Per-profile validation across redshift}

The single-profile multi-$z$ hold-out tests below provide a per-profile
breakdown of emulator accuracy at every trained snapshot in
$z\in[0,0.5]$.

\begin{figure}[htbp]
    \centering
    \includegraphics[width=0.49\linewidth]{CGD_multiz_valid.png}
    \includegraphics[width=0.49\linewidth]{CGED_multiz_valid.png}\\
    \includegraphics[width=0.49\linewidth]{CEP_multiz_valid.png}
    \includegraphics[width=0.49\linewidth]{CEEP_multiz_valid.png}\\
    \includegraphics[width=0.49\linewidth]{CMP_multiz_valid.png}
    \includegraphics[width=0.49\linewidth]{CPP_multiz_valid.png}\\
    \includegraphics[width=0.49\linewidth]{CTP_multiz_valid.png}
    \includegraphics[width=0.49\linewidth]{CYP_multiz_valid.png}
    \caption{Per-profile multi-$z$ hold-out validation. Each subplot
    shows the emulator prediction (dashed) versus three test sims
    (solid) at every trained snapshot ($z=0$ to $z\approx 0.5$). The
    most challenging profiles are the metallicity (CMP, large dynamic
    range) and electron entropy (CEEP, steep gradient at large $r$);
    the others are emulated to better than $\sim$10\% across the
    radial range.}
    \label{fig:profiles_multiz_valid_grid}
\end{figure}

\FloatBarrier

\section{Parameter sensitivity}

\begin{figure}[htbp]
    \centering
    \includegraphics[width=1\linewidth]{profiles_all_sensitivity_z0.png}
    \caption{Per-parameter sensitivity at $z=0$. Rows index profile
    type; columns index parameter, ordered cosmology--first
    ($\omega_{\rm m}$, $\sigma_8$) then subgrid
    ($\kappa_{\rm w}, e_{\rm w}, M_{\rm seed}, v_{\rm kin}, \epsilon_{\rm kin}$).
    Within each panel the emulator is swept across the design range of
    that parameter while the others are held at the design mean.
    Dotted vertical lines mark the rightmost radius where the
    fractional spread falls below 5\%, i.e.\ the outer-radius below
    which the parameter has no measurable impact. Cluster cores show
    the strongest response to subgrid feedback; outer profiles
    ($r\gtrsim R_{500c}$) are nearly degenerate.}
    \label{fig:profiles_sens_z0}
\end{figure}

\begin{figure}[htbp]
    \centering
    \includegraphics[width=1\linewidth]{profiles_all_z_sensitivity.png}
    \caption{Redshift sensitivity at fiducial subgrid+cosmology
    parameters. Each panel shows the emulator prediction over
    $z\in[0,0.5]$ for one profile. Gas density, pressure and Compton-$y$
    show the strongest redshift evolution; metallicity evolves
    smoothly with small amplitude across this range.}
    \label{fig:profiles_z_sens}
\end{figure}

\subsection{Cluster-gas-density per-parameter sensitivity}

\begin{figure}[htbp]
    \centering
    \includegraphics[width=0.6\linewidth]{CGD_sensi_HACC5p.png}
    \caption{CGD sensitivity to each of the seven parameters. The
    strongest impact is on the inner profile from kinetic feedback
    ($v_{\rm kin}, \epsilon_{\rm kin}$) and on the overall amplitude
    from $\sigma_8$ and $\omega_{\rm m}$.}
    \label{fig:cgd_sensi}
\end{figure}

\FloatBarrier

\section{Redshift sweep at fiducial parameters}

\begin{figure}[htbp]
    \centering
    \includegraphics[width=1\linewidth]{profiles_all_z_sweep_sim0.png}
    \caption{Emulator redshift sweep, holding the parameters fixed at
    one test simulation. Each row is a profile, each column a redshift
    snapshot. The colour scale follows redshift.}
    \label{fig:profiles_z_sweep_sim0}
\end{figure}

\begin{figure}[htbp]
    \centering
    \includegraphics[width=0.95\linewidth]{profiles_multiz_redshift_interp.png}
    \caption{Redshift interpolation diagnostic for all profiles at one
    test sim: the emulator's continuous-$z$ prediction is overlaid for
    five redshifts spanning the trained range. The transitions between
    snapshots are smooth and monotonic in the regions where the
    underlying simulation is monotonic, as required for the calibration
    likelihood to be well-conditioned.}
    \label{fig:profiles_z_interp}
\end{figure}

\subsection{Single-profile $z$ sweeps}

\begin{figure}[htbp]
    \centering
    \includegraphics[width=0.95\linewidth]{CGD_z_sweep.png}\\[1ex]
    \includegraphics[width=0.95\linewidth]{CGED_z_sweep.png}
    \caption{Redshift evolution of the emulated CGD (top) and CGED
    (bottom) at the design points. The amplitude in the core
    and at intermediate radii decreases with $z$, while the outer
    profile stays approximately self-similar.}
    \label{fig:cgd_cged_z_sweep}
\end{figure}

\FloatBarrier

\section{Comparison with published cluster-profile observations}

\begin{figure}[htbp]
    \centering
    \includegraphics[width=1\linewidth]{profiles_emu_vs_obs_z0.png}
    \caption{Emulator predictions at $z=0$ (coloured curves, one per
    held-out test sim, with $90\%$ GP bands) overlaid against external
    cluster-profile observations (black markers). Observations include
    McDonald+2017 (CGD/CGED, $z=0$ bin), Pakmor+2022 ($M=10^{14}\,M_\odot$
    bin), Lovisari+2019 metallicity (relaxed sample) and Braspenning+2023
    (CC sample). CEEP and CYP have no external dataset and are omitted
    from this figure. The emulator predictions sit within the observed
    band over most of the radial range, with mild tension at small radii
    for the metallicity and at the outermost bins of the temperature
    profile.}
    \label{fig:profiles_emu_vs_obs}
\end{figure}

\subsection{CGD-only training-set view}

\begin{figure}[htbp]
    \centering
    \includegraphics[width=0.85\linewidth]{CGD_design.png}\\[1ex]
    \includegraphics[width=0.85\linewidth]{CGED_design.png}
    \caption{Training ensembles for the cluster gas density (top) and
    electron density (bottom), coloured by a subgrid parameter. Black
    markers show the published McDonald+2017 datapoints used for
    calibration in the inference pipeline. The mass/radius cut applied
    to the emulator is shaded.}
    \label{fig:cgd_cged_design}
\end{figure}

\begin{figure}[htbp]
    \centering
    \includegraphics[width=0.55\linewidth]{CGD_valid.png}
    \caption{CGD hold-out validation at $z=0$ for three test
    simulations.}
    \label{fig:cgd_valid}
\end{figure}

\end{document}
